\subsection{Resultants}
In view of \eqref{eq:stress-resultants}, the warping shapes \eqref{eq:def-eight-shapes} introduce the following generalized resultants:

\begingroup
\renewcommand{\arraystretch}{1.5}
\setlength{\extrarowheight}{.6ex}
\[
\begin{array}{@{}l >{\displaystyle}l @{\qquad\quad} l >{\displaystyle}l @{}}
\multicolumn{2}{@{}l}{\text{\textbf{Generalized bi-shears.}}}
 & \multicolumn{2}{l@{}}{\text{\textbf{Generalized bi-moments.}}} \\[0.25ex]
\bm{v}_{(b)} & \text{— workless}
 & \bm{w}_{(b)} & \triangleq \int_{\Section} \PoissonFlexure^{\!\top}\,\bm{\tau}\dA \\
\bm{v}_{(a)} & \text{— workless}
 & \bm{w}_{(a)} & \triangleq \int_{\Section} \bm{W}_{(a)}^{\!\top}\,\bm{\tau}\dA \\
\bm{v}_{(\varepsilon)} & \text{— workless}
 & \bm{w}_{(\varepsilon)} & \triangleq \int_{\Section} \bm{W}_{(\varepsilon)}^{\!\top}\,\bm{\tau}\dA \\
\bm{v}_{(\WarpTwistIesan)} & \triangleq \int_{\Section} (\nabla\bm\WarpTwistIesan)\,\bm\tau\dA
 & \bm{w}_{(\psi)} & \triangleq \int_{\Section} \WarpTwistIesan\,\sigma\dA \\
\bm{v}_{(\WarpShearIesan)} & \triangleq \int_{\Section} (\nabla\bm\WarpShearIesan)\,\bm\tau\dA
 & \bm{w}_{(\psi)} & \triangleq \int_{\Section} \bm{\WarpShearIesan}\,\sigma\dA \\
\end{array}
\]
\endgroup

% \paragraph{Classical shear resultant.}
Define $\ShearForce \triangleq \int_{\Section}\bm\tau\dA = \bm{v}_{(r)}$. 
A tensor form is
\begin{equation}
\ShearForce = \BishearTensorRB\,\ShiftB\bm\gamma,
\qquad
\BishearTensorRB \coloneq \int_{\Section} G\,\mathbf{B}_{(b)}\dA,
\label{eq:bishear-r-tensor}
\end{equation}
where \([\BishearTensorRB] = \text{FL}^2\).
\textbf{TODO (proof,\cite[Prop. 2.1]{romano2012torsion})}:
\begin{equation}
\BishearTensorRB \equiv \int_{\Section} E\,(\bm{r}-\CentroidLocation)\otimes(\bm{r}-\CentroidLocation)\dA
=: E\,\IntRGoRG,
\label{eq:HF-eq-EJG}
\end{equation}
where recall \(\IntRGoRG\) is defined by \eqref{eq:def-intRGoRG}. 
% \paragraph{Flexural Longitudinal Warping  $\bm{v}_{(\WarpShearIesan)}$ and $\bm{w}_{\WarpShearIesan}$.}
The bi-shear \(\bm{v}_{\WarpShearIesan}\) is an in-plane force obtained by weighting $\bm\tau$ with $\nabla\bm\WarpShearIesan$. 

\begin{equation}
\bm{v}_{(\WarpShearIesan)} \;=\; \BishearTensorPB\,\ShiftB\bm\gamma,
\qquad
\BishearTensorPB \coloneq \int_{\Section} \nabla\bm\WarpShearIesan \,G\mathbf{B}_{(b)} \dA.
\label{eq:bishear-psi-tensor}
\end{equation}
%
Using equilibrium with the Saint-Venant field \eqref{eq:sv-equilibrium} gives
$\bm{v}_{(\WarpShearIesan)}=\int_{\Section}\bm\WarpShearIesan\,\sigma'\dA$.
For the $\bm{r}$-affine Saint-Venant stress
$\sigma(x,\bm{r})=a_{\varepsilon}(\reflenx)+\bm a_\kappa(\reflenx)\cdot(\bm{r}-\CentroidLocation)$
and the gauge $\int_{\Section}\bm\WarpShearIesan\dA=\mathbf{0}$,

\textbf{TODO}: Prove:
\begin{equation*}
  \BishearTensorPB \coloneq \int\nabla\bm{\WarpShearIesan}\,G\,\mathbf{B}_{(b)}\dA
  \equiv \int_{\Section}E \bm\WarpShearIesan \otimes(\bm{r}-\CentroidLocation)\dA.
\end{equation*}
so \(\bm{v}_{(\bm\WarpShearIesan)}\) is driven by the gradients of axial strain and curvatures, via the ``warping moments'' \(\BishearTensorPB\).
\begin{remark}
\emph{Definiteness.}
If $\nu=0$ (homogeneous Neumann data for $\bm\WarpShearIesan\cdot\bm b$) and $\int_{\Section}\bm\WarpShearIesan\dA=\mathbf{0}$, then $\BishearTensorPB$ is symmetric positive definite.
For $\nu\neq 0$, the Poisson coupling produces an inhomogeneous boundary flux; no general symmetry/positivity guarantee holds and the property must be checked for the specific section and boundary data.
\end{remark}

% \paragraph{Flexural Poisson Warping  $\bm{w}_{(b)}$.}
The bi-moment $\bm{w}_{(b)}$ is an in-plane force collecting the boundary-driven, Poisson-coupled part of the shear–warping interaction. 
\begin{equation}
\BishearTensorBB \coloneq \SectionIntegral{\PoissonFlexure^{\!\top}\,G\,\mathbf{B}_{(b)}},
\label{eq:bishear-b-tensor}
\end{equation}
When $\nu=0$, $\PoissonFlexure\equiv\mathbf{0}$ and hence $\bm{w}_{(b)}=\mathbf{0}$. 
If $\bm\tau'=\mathbf{0}$ (Saint-Venant shear independent of $x$), then $\bm{w}_{(b)}'=\mathbf{0}$.

\paragraph{Summary.}
The additional in-plane forces $\bm{v}_{(\WarpShearIesan)}$ and $\bm{w}_{(b)}$ enter the reduced 1D balance/constitutive equations when the cross-section distorts:
\begin{itemize}
  \item $\bm{v}_{(\WarpShearIesan)}$ is a bulk coupling that vanishes only for $\reflenx$-uniform axial stress; otherwise it captures axial-stress-gradient forcing on the shear–warping mode via \eqref{eq:bishear-psi-tensor}.
  \item $\bm{w}_{(b)}$ is a boundary/Poisson coupling that persists even under Saint-Venant assumptions; for $\nu=0$ it vanishes, while for $\nu\neq 0$ it influences effective shear stiffness, shear-center shifts, and flexure–shear coupling.
\end{itemize}


% \clearpage
% \input{40_Resultants-1}
% \clearpage